\section{\lr{attention rollout} و \lr{attention flow}}
\label{sec:rollout-flow}

\lr{attention rollout} و \lr{attention flow} \cite{abnar2020} هر دو یک پرسش را پاسخ می‌دهند: یک عدد خروجی از کدام ورودی‌ها آمده است، وقتی اطلاعات در هر لایه فقط از راه توجه و \lr{residual connection} بین موقعیت‌ها جابه‌جا می‌شود. هر دو روی همان ماتریس‌های ثبت‌شده کار می‌کنند و فقط در قاعده‌ی ترکیب لایه‌ها فرق دارند: \lr{rollout} در طول هر مسیر وزن‌ها را ضرب و مسیرها را جمع می‌کند، و \lr{attention flow} وزن‌ها را ظرفیت لوله می‌گیرد و بیشینه‌ی جریانی را که از یک ورودی به خروجی می‌رسد حساب می‌کند. قاعده‌ی «یک مخلوط‌کننده» در \lr{v2} (بخش~\ref{sec:v2}) شرط لازم هر دو است: چون اطلاعات جز از راه توجه (و \lr{stem} معلوم \lr{K}) بین موقعیت‌ها جابه‌جا نمی‌شود، این گراف کل مسیر اطلاعات را می‌پوشاند.

\subsection{نمادگذاری و \lr{layered graph}}

برای لایه‌ی $l$ از $L$ لایه‌ی توجه، \lr{recorder} ماتریس $A^{h}_l\in\mathbb{R}^{n\times n}$ را برای هر سر $h=1,\dots,H$ پس از \lr{softmax} برمی‌دارد؛ سطر $i$ گیرنده (پرسش) و ستون $j$ فرستنده است و هر سطر جمع ۱ دارد. میانگین سرها و \lr{residual connection}:
\begin{equation}
\overline{A}_l=\frac1H\sum_{h=1}^{H}A^{h}_l,\qquad \hat A_l=\tfrac12\,\overline{A}_l+\tfrac12 I .
\label{eq:ahat}
\end{equation}
وزن $\tfrac12$ برای \lr{residual connection} یک قرارداد است: بلوک $x+\mathrm{Attn}(\mathrm{LN}(x))$ نیمی از اطلاعات را از مسیر توجه و نیمی را از مسیر مستقیم فرض می‌کند. \lr{layered graph} $L+1$ ردیف گره دارد (ردیف ۰ ورودی‌ها، ردیف $L$ خروجی‌ها) و از گره $j$ در ردیف $l-1$ به گره $i$ در ردیف $l$ یک یال با وزن $\hat A_l[i,j]$ می‌رود. یال‌های $j\to j$ مسیر اتصال باقی‌مانده‌اند و بقیه مسیر توجه (شکل~\ref{fig:layered}).

\begin{figure}[htbp]
\centering
% Layered graph of rollout and attention flow.
\begin{latin}
\begin{tikzpicture}[x=2.6cm, y=1.9cm, font=\footnotesize,
  nd/.style={circle, draw, minimum size=7mm, inner sep=0pt, font=\small, fill=white},
  att/.style={-{Stealth[length=1.6mm]}, black!35},
  res/.style={-{Stealth[length=2mm]}, line width=1.2pt, black!70},
  hot/.style={-{Stealth[length=2.2mm]}, line width=1.6pt, black}]
\foreach \l in {0,1,2,3}{\foreach \i in {1,2,3,4}{\node[nd] (n\l\i) at (\i,\l) {$\i$};}}
\foreach \l/\p in {1/0,2/1,3/2}{
  \foreach \i in {1,2,3,4}{\foreach \j in {1,2,3,4}{\ifnum\i=\j\else\draw[att] (n\p\j)--(n\l\i);\fi}}
  \foreach \i in {1,2,3,4}{\draw[res] (n\p\i)--(n\l\i);}
}
\draw[hot] (n01)--(n12); \draw[hot] (n12)--(n22); \draw[hot] (n22)--(n33);
\node[nd, line width=1.4pt] at (3,3) {$3$};
\foreach \l/\t in {0/{inputs (tokens or electrodes)},1/layer 1,2/layer 2,3/{output (layer 3)}}{\node[lab, anchor=west] at (4.35,\l) {\t};}
\node[lab] at (2.5,3.6) {target node: token $t=3$ of the output};
\draw[res] (0.6,-0.7)--(1.0,-0.7); \node[lab, anchor=west] at (1.05,-0.7) {residual edge, capacity $\tfrac12+\tfrac12\overline{A}_{ii}$};
\draw[att] (0.6,-1.05)--(1.0,-1.05); \node[lab, anchor=west] at (1.05,-1.05) {attention edge, capacity $\tfrac12\overline{A}_{ij}$};
\draw[hot] (0.6,-1.4)--(1.0,-1.4); \node[lab, anchor=west] at (1.05,-1.4) {one path from input 1 to the target};
\end{tikzpicture}
\end{latin}

\caption{\lr{layered graph} \lr{rollout} و \lr{attention flow}. وزن هر یال همان $\hat A_l[i,j]$ است. \lr{rollout} روی همه‌ی مسیرها از ورودی به هدف حاصل‌ضرب وزن‌ها را جمع می‌زند؛ \lr{attention flow} وزن‌ها را ظرفیت لوله می‌داند و بیشینه‌ی جریانی را که از یک ورودی به هدف می‌رسد حساب می‌کند.}
\label{fig:layered}
\end{figure}

\subsection{\lr{attention rollout}}

\subsubsection*{تعریف}
\lr{rollout} حاصل‌ضرب ماتریس‌های لایه‌ها از آخر به اول است، و این حاصل‌ضرب دقیقاً جمع روی همه‌ی مسیرهای \lr{layered graph} است:
\begin{equation}
R=\hat A_L\hat A_{L-1}\cdots\hat A_1,\qquad
R[t,\tau]=\sum_{\tau=i_0\to i_1\to\cdots\to i_L=t}\ \prod_{l=1}^{L}\hat A_l[i_l,i_{l-1}] .
\label{eq:rollout}
\end{equation}
هر مسیر به اندازه‌ی حاصل‌ضرب وزن‌هایش سهم می‌گیرد، پس یک یال ضعیف کل مسیر را تضعیف می‌کند. چون همه‌ی $\hat A_l$ سطرهایی با جمع ۱ و درایه‌های نامنفی دارند، حاصل‌ضرب آن‌ها هم همین‌طور است:
\begin{equation}
\hat A_l\mathbf 1=\mathbf 1\ \ \forall l\quad\Longrightarrow\quad R\,\mathbf 1=\hat A_L\cdots\hat A_1\mathbf 1=\mathbf 1,\qquad \sum_\tau R[t,\tau]=1 ,
\end{equation}
و سطر $t$ از $R$ یک توزیع روی ورودی‌هاست: سهم هر ورودی در خروجی \lr{token} $t$.

\subsubsection*{\lr{rollout} در هر مسیر مدل}
همین حاصل‌ضرب، مرحله به مرحله، برای هر خروجی:
\begin{itemize}
\item \textbf{\lr{K}، زمان} (سه لایه‌ی زمانی سینماتیک و \lr{stem}): $R_K=\hat A^{(k)}_3\hat A^{(k)}_2\hat A^{(k)}_1\,S$
\item \textbf{\lr{E}، زمان} (دو لایه‌ی زمانی \lr{EEG}): $R_E=\hat A^{(t)}_2\hat A^{(t)}_1$
\item \textbf{\lr{E}، الکترودها در \lr{token} $\tau$} (دو لایه‌ی گراف الکترود و \lr{pooling}): $e_\tau=\pi_\tau\,\hat A^{(e)}_2(\tau)\,\hat A^{(e)}_1(\tau)\in\mathbb{R}^{32}$
\item \textbf{\lr{E}، زمان در الکترود}: $M_E[\tau,e]=R_E[t,\tau]\;e_\tau[e]$
\item \textbf{$F_{eeg}$}: $R_F=c_t\,R_E$، که $c_t$ سطر $t$ \lr{cross-attention} است (بدون \lr{residual connection})
\item \textbf{\lr{muscle graph}} (دو لایه): $R_M=\hat A^{(m)}_2\hat A^{(m)}_1\in\mathbb{R}^{5\times5}$
\end{itemize}
$S$ ردپای \lr{convolutional stem} \lr{K} است: ماتریسی نواری با پهنای ۵ که وزن هر جابه‌جایی برابر انرژی هسته در آن جابه‌جایی است و سطرهایش جمع ۱ دارند. $\pi_\tau$ وزن‌های \lr{pooling} الکترودها در \lr{token} $\tau$ است (جمع ۱). چون همه‌ی عامل‌ها جمع سطری ۱ دارند، $\sum_{\tau,e}M_E[\tau,e]=1$ هم برقرار است.

\subsubsection*{نمایش}
\lr{token} خود رخداد و همسایه‌هایش (ردپای ۵ توکنی \lr{stem}) همیشه بیشترین سهم را از مسیر مستقیم می‌گیرند و بقیه‌ی نقشه را پنهان می‌کنند. شکل‌ها سهم آن‌ها را جداگانه می‌نویسند و منحنی سهم بقیه‌ی \lr{tokens} را نشان می‌دهند:
\begin{equation}
\mathrm{own}=\sum_{\tau\in\mathcal{O}}R[t,\tau],\qquad s(\tau)=\frac{R[t,\tau]}{1-\mathrm{own}}\quad(\tau\notin\mathcal{O}),
\end{equation}
که $\mathcal{O}$ پنج \lr{token} اطراف $t$ است. نقشه‌ی هر رخداد جداگانه ساخته و سپس روی رخدادها میانگین و \lr{SEM} گرفته می‌شود.

\subsubsection*{نسخه‌ی وزن‌دار با گرادیان}
در \lr{gradient-weighted rollout} (بخش~\ref{sec:interp}) فقط $\overline{A}_l$ در رابطه‌ی~\eqref{eq:ahat} با $\tilde A_l=\mathrm{rownorm}\big(\frac1H\sum_h(A^{h}_l\odot\partial y/\partial A^{h}_l)^{+}\big)$ عوض می‌شود و بقیه‌ی مراحل همین است؛ نتیجه برای هر عضله جداست.

\subsubsection*{محدودیت‌ها}
\lr{rollout} فقط وزن‌های توجه را می‌بیند و آنچه بردارهای مقدار حمل می‌کنند در آن نیست. لایه‌های بدون توجه (\lr{embedding}، \lr{feed-forward}، \lr{decoder}) در آن وارد نمی‌شوند، و وزن \lr{۰٫۵} \lr{residual connection} انتخاب ماست، نه چیزی که مدل آموخته باشد.

\begin{figure}[p]
\centering
% Attention rollout: the whole process in the style of a paper figure.
\begin{latin}\resizebox{\textwidth}{!}{%
\begin{tikzpicture}[font=\footnotesize,
  stepc/.style={circle, fill=black!80, text=white, font=\scriptsize\bfseries, inner sep=1.2pt, minimum size=4.6mm},
  stitle/.style={lab, font=\small\bfseries, anchor=west}]
% ---------------- step 1: record
\node[stepc] at (-0.4,0.3) {1};
\node[stitle] at (-0.05,0.3) {Record every attention layer in one forward pass (after softmax, per head)};
\foreach \l/\x/\s in {1/1.6/0,2/6.1/-1,3/10.6/1}{
  \foreach \h in {2,1,0}{\attnmap{r1\l\h}{\x+0.18*\h,-1.65+0.18*\h}{8}{0.2}{\s+\h-1}{intc}{0}}
  \node[lab] at (\x+0.2,-2.75) {layer \l: $A^{h}_{\l}$, heads $h=1..H$};}
\node[lab, align=left, anchor=west] at (12.75,-1.55) {row $i$ = receiving\\ token (query),\\ column $j$ = source};

% ---------------- step 2: heads and residual
\node[stepc] at (-0.4,-3.45) {2};
\node[stitle] at (-0.05,-3.45) {Average the heads and add the residual path};
\foreach \l/\x/\s in {1/1.6/0,2/6.1/-1,3/10.6/1}{
  \attnmap{r2\l}{\x,-5.05}{8}{0.2}{\s}{intc}{30}
  \node[lab] at (\x,-6.1) {$\hat A_{\l}=\tfrac12\,\overline{A}_{\l}+\tfrac12 I$};
  }
\node[lab, align=left, anchor=west] at (12.75,-5.05) {$\overline{A}_l=\frac1H\sum_h A^{h}_l$\\ rows still sum to 1\\ $I$: the skip connection};

% ---------------- step 3: multiply
\node[stepc] at (-0.4,-6.85) {3};
\node[stitle] at (-0.05,-6.85) {Multiply the layers from the last to the first};
\attnmap{m3}{1.2,-8.4}{8}{0.2}{1}{intc}{30}\node[lab] at (1.2,-9.45) {$\hat A_3$};
\node at (2.4,-8.4) {$\cdot$};
\attnmap{m2}{3.6,-8.4}{8}{0.2}{-1}{intc}{30}\node[lab] at (3.6,-9.45) {$\hat A_2$};
\node at (4.8,-8.4) {$\cdot$};
\attnmap{m1}{6.0,-8.4}{8}{0.2}{0}{intc}{30}\node[lab] at (6.0,-9.45) {$\hat A_1$};
\node at (7.2,-8.4) {$\cdot$};
\attnmap{ms}{8.4,-8.4}{8}{0.2}{0}{kinc}{25}\node[lab] at (8.4,-9.45) {$S$ (K only: stem)};
\node at (9.6,-8.4) {$=$};
\attnmap{rr}{10.9,-8.4}{8}{0.2}{0}{black}{10}\node[lab] at (10.9,-9.45) {$R$};
\draw[red!75!black, line width=1pt] (10.1,-8.4+0.0) rectangle (11.7,-8.2);
\node[lab, text=red!75!black, anchor=west] at (11.75,-8.3) {row $t$ = the explained token};
\node[lab, align=left, anchor=west] at (12.75,-9.05) {$R[t,\tau]=\sum_{\text{paths }\tau\to t}\prod_l \hat A_l$\\ $\sum_\tau R[t,\tau]=1$};

% ---------------- step 4: read the row
\node[stepc] at (-0.4,-10.2) {4};
\node[stitle] at (-0.05,-10.2) {Read row $t$: the share of every input in the explained output};
\draw[black!60] (0.3,-12.7)--(8.3,-12.7);
\foreach \v [count=\k] in {0.10,0.12,0.15,0.2,0.28,0.42,0.6,0.5,0.35,1.55,0.3,0.22,0.16,0.12,0.1,0.09,0.08,0.08,0.07,0.07}{
  \fill[intc!70] (0.3+0.38*\k,-12.7) rectangle ++(0.28,\v);}
\fill[black!70] (0.3+0.38*10,-12.7) rectangle ++(0.28,1.55);
\node[lab, anchor=west] at (4.35,-11.0) {own token (skip path) $\to$ written apart};
\node[lab] at (4.3,-13.0) {input token $\tau$ (time)};

% ---------------- step 5: what R means in each stream
\node[stepc] at (-0.4,-13.6) {5};
\node[stitle] at (-0.05,-13.6) {The same product, stage by stage, for every output};
\node[kbox, text width=3.6cm, align=left] (b1) at (1.75,-15.05) {\textbf{K}: $R_K=\hat A_3\hat A_2\hat A_1 S$\\ time map, 400 tokens};
\node[ebox, text width=3.6cm, align=left] (b2) at (5.95,-15.05) {\textbf{E}, time: $R_E=\hat A_2\hat A_1$\\ (EEG temporal layers)};
\node[ebox, text width=4.4cm, align=left] (b3) at (10.45,-15.05) {\textbf{E}, electrodes at token $\tau$:\\ \mbox{$e_\tau=\pi_\tau\,\hat A^{(e)}_2\hat A^{(e)}_1$}, $\pi$ = pooling row};
\headgraph{hr}{14.2,-15.05}{0.75}
\node[ebox, text width=5.6cm, align=left] (b4) at (3.0,-16.75) {\textbf{E}, time $\times$ electrode:\\ $M_E[\tau,e]=R_E[t,\tau]\;e_\tau[e]$};
\node[fbox, text width=4.2cm, align=left] (b5) at (8.55,-16.75) {\textbf{F$_{eeg}$}: cross-attention $C$ last,\\ no residual: $C\,R_E$};
\node[box, fill=intc!8, text width=4.4cm, align=left] (b6) at (13.45,-16.75) {\textbf{muscle graph}: $\hat A^{(m)}_2\hat A^{(m)}_1$\\ 5$\times$5, row = receiving muscle};

% ---------------- step 6: average over events
\node[stepc] at (-0.4,-17.95) {6};
\node[stitle] at (-0.05,-17.95) {One map per lift event, then mean $\pm$ SEM over events (never a single window)};
\fill[intc!20] plot[smooth, domain=0.3:8.3, samples=50] (\x,{-19.6+0.55*exp(-((\x-3.6)^2)/0.8)+0.25*exp(-((\x-6.9)^2)/0.5)+0.12}) -- (8.3,-19.75) -- plot[smooth, domain=8.3:0.3, samples=50] (\x,{-19.75+0.55*exp(-((\x-3.6)^2)/0.8)+0.25*exp(-((\x-6.9)^2)/0.5)+0.05}) -- cycle;
\draw[intc!80!black, line width=0.9pt] plot[smooth, domain=0.3:8.3, samples=50] (\x,{-19.68+0.55*exp(-((\x-3.6)^2)/0.8)+0.25*exp(-((\x-6.9)^2)/0.5)+0.08});
\draw[black!60] (0.3,-19.75)--(8.3,-19.75); \draw[black!50, dashed] (4.3,-19.8)--(4.3,-18.75);
\node[lab] at (4.3,-20.05) {time from the event (s); dashed = event};
\node[box, fill=black!3, text width=7.4cm, align=left] at (12.6,-19.35) {\textbf{gradient-weighted variant}: replace $\overline{A}_l$ by\\ $\tilde A_l=\mathrm{rownorm}\Big(\frac1H\sum_h\big(A^{h}_l\odot\frac{\partial y}{\partial A^{h}_l}\big)^{+}\Big)$\\ and run steps 2--6 unchanged: one map per muscle.};
\end{tikzpicture}}\end{latin}

\caption{فرایند کامل \lr{attention rollout}. (۱) \lr{recorder} توجه هر لایه و هر سر را در یک گذر برمی‌دارد. (۲) سرها میانگین گرفته و \lr{residual connection} به‌صورت $\tfrac12 I$ اضافه می‌شود. (۳) ماتریس‌ها از لایه‌ی آخر به اول ضرب می‌شوند (برای \lr{K} ردپای \lr{stem} $S$ هم). (۴) سطر \lr{token} توضیح‌داده‌شده سهم هر ورودی است؛ سهم \lr{token} خود رخداد جداگانه نوشته می‌شود. (۵) همین حاصل‌ضرب برای زمان، الکترودها، زمان در الکترود، اصلاح \lr{EEG} و \lr{muscle graph}. (۶) نقشه‌ی هر رخداد ساخته و روی رخدادها میانگین گرفته می‌شود. کادر پایین: نسخه‌ی وزن‌دار با گرادیان.}
\label{fig:rollout-process}
\end{figure}

\subsection{\lr{attention flow}}

\subsubsection*{تعریف}
در \lr{attention flow} هر یال \lr{layered graph} یک لوله با ظرفیت $c_e=\hat A_l[i,j]$ است. برای ورودی (منبع) $\tau$ و گره هدف $t$، جریان یک تابع $f$ روی یال‌هاست و بیشینه‌ی مقدار آن جواب است:
\begin{equation}
\mathrm{flow}(\tau\to t)=\max_{f}\ \sum_{e\,\in\,\mathrm{out}(\tau)} f_e
\quad\text{که}\quad 0\le f_e\le c_e,\qquad \sum_{e\,\in\,\mathrm{in}(v)}f_e=\sum_{e\,\in\,\mathrm{out}(v)}f_e\ \ \text{برای هر گره میانی }v .
\label{eq:flow}
\end{equation}
طبق قضیه‌ی \lr{maximum flow} و \lr{minimum cut}، این مقدار برابر ظرفیت کم‌هزینه‌ترین برشی است که $\tau$ را از $t$ جدا می‌کند:
\begin{equation}
\mathrm{flow}(\tau\to t)=\min_{(\mathcal{S},\mathcal{T}):\ \tau\in\mathcal{S},\ t\in\mathcal{T}}\ \sum_{e:\,\mathcal{S}\to\mathcal{T}}c_e .
\end{equation}
فرق با \lr{rollout}: \lr{rollout} در طول هر مسیر ضرب می‌کند، اما جریان مسیرهای موازی را جمع می‌کند و هر مسیر فقط با باریک‌ترین لوله‌اش محدود می‌شود. به همین دلیل جریان هیچ‌گاه از سهم \lr{rollout} کمتر نیست و وقتی توجه پخش است، زود به سقف می‌رسد.

\subsubsection*{بازه‌بندی زمان}
حل رابطه‌ی~\eqref{eq:flow} روی ۴۰۰ \lr{tokens} برای هر منبع بسیار پرهزینه است. زمان به بازه‌های ۴ توکنی (۸۰ میلی‌ثانیه) تقسیم می‌شود: سطرهای گیرنده‌ی هر بازه میانگین و ستون‌های فرستنده جمع می‌شوند،
\begin{equation}
\hat A^{bin}_l[a,b]=\frac{1}{|a|}\sum_{i\in a}\sum_{j\in b}\hat A_l[i,j],
\end{equation}
که سطرها را همچنان با جمع ۱ نگه می‌دارد. برای \lr{K} ماتریس \lr{stem} $S$ هم با همین قاعده بازه‌بندی و پیش از لایه‌ی اول گذاشته می‌شود؛ برای $F_{eeg}$ سطر \lr{cross-attention} لایه‌ی آخر است.

\subsubsection*{جریان الکترودها}
برای \lr{E} و $F_{eeg}$، در هر بازه‌ی زمانی $\tau$ گراف دیگری هست: دو لایه‌ی \lr{electrode graph} ($32\to32\to32$، با \lr{residual connection}) و ردیف \lr{pooling} ($32\to1$ با ظرفیت $\pi_\tau[e]$ و بدون \lr{residual connection}). $\mathrm{flow}_{elec}(\tau,e)$ \lr{maximum flow} از الکترود $e$ تا \lr{token} پول‌شده است. چون \lr{token} پول‌شده تنها راه خروج اطلاعات الکترودهای یک گام زمانی است، جریان یک الکترود در یک زمان تا خروجی کمینه‌ی دو مرحله است:
\begin{equation}
\mathrm{flow}(\tau,e)=\min\big(\mathrm{flow}_{time}(\tau),\ \mathrm{flow}_{elec}(\tau,e)\big).
\end{equation}

\subsubsection*{خط پایه‌ی توجه یکنواخت}
همان محاسبه برای مدلی تکرار می‌شود که در هر لایه به همه‌ی $n$ موقعیت برابر توجه می‌کند، با همان \lr{residual connection} و همان \lr{stem}:
\begin{equation}
\hat A^{u}=\tfrac12\cdot\tfrac1n J+\tfrac12 I,\qquad \mathrm{excess}(\tau)=\mathrm{flow}(\tau)-\mathrm{flow}^{u}(\tau),
\end{equation}
که $J$ ماتریس تمام‌یک است. فقط $\mathrm{excess}$ به الگوی توجه آموخته‌شده مربوط است؛ $\mathrm{flow}^{u}$ چیزی است که حتی یک مدل آموزش‌ندیده هم می‌دهد.

\begin{figure}[p]
\centering
% Attention flow: the whole process in the style of a paper figure.
\begin{latin}\resizebox{\textwidth}{!}{%
\begin{tikzpicture}[font=\footnotesize,
  stepc/.style={circle, fill=black!80, text=white, font=\scriptsize\bfseries, inner sep=1.2pt, minimum size=4.6mm},
  stitle/.style={lab, font=\small\bfseries, anchor=west},
  fn/.style={circle, draw=black!60, fill=white, minimum size=5.5mm, inner sep=0pt, font=\scriptsize},
  capedge/.style={black!30, line width=0.5pt},
  resedge/.style={black!65, line width=1.1pt},
  fl/.style={-{Stealth[length=1.8mm]}, fusc!70!black, line width=1.4pt}]
% ---------------- step 1: same matrices, binned
\node[stepc] at (-0.4,0.3) {1};
\node[stitle] at (-0.05,0.3) {Start from the same layer matrices $\hat A_l=\tfrac12\overline{A}_l+\tfrac12I$ and bin the time axis};
\attnmap{f1}{1.6,-1.35}{8}{0.2}{0}{intc}{30}
\node[lab] at (1.6,-2.4) {$\hat A_l$: 400 $\times$ 400 tokens};
\draw[arr] (2.6,-1.35)--(4.1,-1.35) node[lab, midway, above]{bins of 4};
\attnmap{f2}{5.1,-1.35}{4}{0.4}{0}{intc}{30}
\node[lab] at (5.1,-2.4) {100 $\times$ 100 bins (80 ms)};
\node[box, fill=black!3, text width=7.0cm, align=left] at (11.2,-1.35) {$\hat A^{bin}[a,b]=\dfrac{1}{|a|}\displaystyle\sum_{i\in a}\sum_{j\in b}\hat A[i,j]$\\[2pt] queries averaged, sources summed: rows still sum to 1, and the graph shrinks from $L\cdot400^2$ to $L\cdot100^2$ edges.};

% ---------------- step 2: the flow network
\node[stepc] at (-0.4,-3.1) {2};
\node[stitle] at (-0.05,-3.1) {Read the layers as a flow network: every edge is a pipe with capacity $\hat A_l[i,j]$};
\foreach \l in {0,1,2,3}{\foreach \i in {1,2,3,4}{\node[fn] (g\l\i) at (0.9+1.3*\i,-8.1+1.3*\l) {};}}
\foreach \l/\p in {1/0,2/1,3/2}{
  \foreach \i in {1,2,3,4}{\foreach \j in {1,2,3,4}{\ifnum\i=\j\else\draw[capedge] (g\p\j)--(g\l\i);\fi}}
  \foreach \i in {1,2,3,4}{\draw[resedge] (g\p\i)--(g\l\i);}}
\node[fn, fill=fusc!45, draw=fusc!70!black] at (g01) {$\tau$};
\node[fn, fill=red!35, draw=red!70!black] at (g33) {$t$};
\draw[fl] (g01)--(g12); \draw[fl] (g12)--(g22); \draw[fl] (g22)--(g33);
\draw[fl] (g01)--(g11); \draw[fl] (g11)--(g23); \draw[fl] (g23)--(g33);
\draw[red!70!black, dashed, line width=1pt] (1.6,-4.85)..controls (3.5,-5.0) and (5.6,-4.5)..(6.4,-4.7);
\node[lab, text=red!70!black, anchor=east] at (1.55,-4.85) {minimum cut};
\node[lab, anchor=east] at (1.9,-8.1) {inputs}; \node[lab, anchor=east] at (1.9,-6.8) {layer 1};
\node[lab, anchor=east] at (1.9,-5.5) {layer 2}; \node[lab, anchor=east] at (1.9,-4.2) {output};
\node[lab, text=fusc!60!black] at (4.15,-8.65) {source: input bin $\tau$ (one problem per source)};
\node[box, fill=black!3, text width=7.0cm, align=left] at (11.2,-6.15) {\textbf{maximum flow} from source $\tau$ to target $t$:\\[3pt]
 $\displaystyle\mathrm{flow}(\tau\!\to\!t)=\max_{f}\sum_{e\,\text{out of}\,\tau} f_e$\\[2pt]
 subject to $0\le f_e\le c_e$ (capacity $c_e=\hat A_l[i,j]$)\\ and inflow = outflow at every inner node.\\[3pt]
 \textbf{max-flow = min-cut}: the value equals the capacity of the cheapest set of pipes that separates $\tau$ from $t$.\\[3pt]
 Rollout multiplies along each path; flow adds parallel paths, each limited by its narrowest pipe.};

% ---------------- step 3: solve
\node[stepc] at (-0.4,-9.5) {3};
\node[stitle] at (-0.05,-9.5) {Solve one maximum-flow problem per source bin (100 per window)};
\node[box, text width=5.6cm, align=left] at (2.8,-10.85) {capacities $\times 10^{6}$, rounded to integers\\ \texttt{scipy...csgraph.maximum\_flow}\\ result divided by $10^{6}$};
\draw[black!60] (6.3,-11.7)--(14.3,-11.7);
\foreach \v [count=\k] in {0.42,0.44,0.43,0.46,0.47,0.45,0.48,0.47,0.5,0.95,0.48,0.46,0.45,0.44,0.46,0.43,0.42,0.41,0.42,0.4}{
  \fill[intc!65] (6.3+0.4*\k-0.3,-11.7) rectangle ++(0.26,\v*1.3);}
\fill[black!70] (6.3+0.4*10-0.3,-11.7) rectangle ++(0.26,0.95*1.3);
\draw[black!55, dashed, line width=0.8pt] (6.3,-11.7+0.5*1.3)--(14.3,-11.7+0.5*1.3) node[lab, anchor=west]{uniform: 0.5};
\node[lab] at (10.3,-12.0) {source bin $\tau$ (time); dark bar = the target's own bin};

% ---------------- step 4: baseline
\node[stepc] at (-0.4,-12.75) {4};
\node[stitle] at (-0.05,-12.75) {Subtract the flow of a model that attends to everything equally};
\node[box, text width=7.3cm, align=left] at (3.6,-13.95) {uniform layers $\tfrac12\,\tfrac1n J+\tfrac12 I$ (same residual, same stem)\\[2pt] $\mathrm{excess}(\tau)=\mathrm{flow}(\tau)-\mathrm{flow}_{uniform}(\tau)$\\[2pt] two uniform layers: every source gets exactly $\tfrac12$};
\draw[black!60] (7.9,-13.5)--(14.3,-13.5) node[lab, anchor=west]{0};
\draw[intc!80!black, line width=0.9pt] plot[smooth, domain=7.9:14.3, samples=40] (\x,{-13.5-0.45-0.18*sin(\x*140)-0.1*sin(\x*310)});
\node[lab, anchor=west, align=left] at (7.9,-14.6) {excess below 0: the learned attention is\\ more concentrated than uniform};

% ---------------- step 5: electrodes
\node[stepc] at (-0.4,-15.6) {5};
\node[stitle] at (-0.05,-15.6) {Electrode flow (E and F$_{eeg}$): per time bin, through both electrode layers into the pooled token};
\foreach \l in {0,1,2}{\foreach \i in {1,...,5}{\node[fn, minimum size=4.2mm, fill=eegc!10, draw=eegc!70!black] (e\l\i) at (0.7+0.75*\i,-19.2+0.95*\l) {};}}
\node[fn, minimum size=5mm, fill=eegc!35, draw=eegc!70!black] (pool) at (2.95,-16.35) {p};
\foreach \l/\p in {1/0,2/1}{\foreach \i in {1,...,5}{\foreach \j in {1,...,5}{\draw[capedge] (e\p\j)--(e\l\i);}}}
\foreach \i in {1,...,5}{\draw[eegc!60, line width=0.7pt] (e2\i)--(pool);}
\node[lab, anchor=west] at (4.95,-19.2) {32 electrodes};
\node[lab, anchor=west] at (4.95,-18.25) {graph layer 1};
\node[lab, anchor=west] at (4.95,-17.3) {graph layer 2};
\node[lab, anchor=west, align=left] at (3.4,-16.35) {pooled token\\ (capacities $\pi_e$, sum 1)};
\node[box, fill=black!3, text width=7.2cm, align=left] at (11.35,-17.85) {$\mathrm{flow}_{elec}(\tau,e)$: maximum flow from electrode $e$ to the pooled token of bin $\tau$.\\[3pt]
 The pooled token is the only exit of a time step, so the input flow is\\ $\mathrm{flow}(\tau,e)=\min\big(\mathrm{flow}_{time}(\tau),\ \mathrm{flow}_{elec}(\tau,e)\big)$.\\[3pt]
 If $\max_e-\min_e<10\,\%$ of the mean, every electrode reaches the same cap: the map carries no spatial information and is not drawn.};
\end{tikzpicture}}\end{latin}

\caption{فرایند کامل \lr{attention flow}. (۱) همان ماتریس‌های لایه‌ها، با بازه‌بندی زمان به ۱۰۰ بازه‌ی ۸۰ میلی‌ثانیه‌ای. (۲) لایه‌ها به‌صورت شبکه‌ی جریان خوانده می‌شوند: هر یال لوله‌ای با ظرفیت $\hat A_l[i,j]$ است؛ \lr{maximum flow} از منبع $\tau$ تا هدف $t$ برابر ظرفیت \lr{minimum cut} است. (۳) برای هر بازه‌ی منبع یک مسئله‌ی \lr{maximum flow} با ظرفیت‌های صحیح حل می‌شود. (۴) جریان یک مدل با توجه یکنواخت کم می‌شود؛ مقدار زیر صفر یعنی توجه آموخته‌شده متمرکزتر از یکنواخت است. (۵) جریان الکترودها از دو لایه‌ی گراف الکترود و \lr{pooling} عبور می‌کند و جریان ورودی کمینه‌ی جریان زمانی و جریان الکترود است.}
\label{fig:flow-process}
\end{figure}

\subsection{پیاده‌سازی}

تابع \lr{trace} مدل را یک بار با \lr{recorder} اجرا می‌کند و برای هر لایه تنسور $(B,H,n,n)$ را برمی‌دارد. \lr{layer\_matrix} رابطه‌ی~\eqref{eq:ahat} را می‌سازد و \lr{rollout} ماتریس‌ها را به ترتیب از لایه‌ی آخر به اول ضرب می‌کند. \lr{rollout} با تابع \lr{rollout} مدل پایه روی ماتریس‌های دست‌ساز آزموده شده است.

برای جریان، \lr{bin\_matrix} بازه‌بندی را انجام می‌دهد و \lr{attention\_flow} برای هر منبع یک مسئله‌ی \lr{maximum flow} با \lr{SciPy} حل می‌کند. این الگوریتم ظرفیت صحیح می‌خواهد، پس ظرفیت‌ها در $10^{6}$ ضرب و گرد می‌شوند و جواب دوباره بر $10^{6}$ تقسیم می‌شود. هر پنجره ۱۰۰ مسئله دارد، پس جریان برای هر خروجی و هر نوع رخداد فقط روی حداکثر ۶ رخداد حساب می‌شود. جریان با \lr{maximum flow} دست‌حساب‌شده روی گلوگاه‌های ساده آزموده شده و هیچ‌گاه کمتر از \lr{rollout} نیست.

\begin{figure}[htbp]
\centering
% Implementation of rollout (left) and attention flow (right) on the time path.
\begin{latin}
\begin{tikzpicture}[node distance=5mm and 5mm, font=\footnotesize, box/.append style={text width=4.2cm}]
\node[io] (x) {one 8 s window, $(1,400,\cdot)$};
\node[box, below=of x] (tr) {\texttt{trace}: one forward pass with the recorder};
\node[box, below=of tr] (A) {$A^{h}_l$ of every layer, $(1,H,n,n)$};
\node[box, below=of A] (L) {\texttt{layer\_matrix}: $\tfrac12\overline{A}_l+\tfrac12I$};
\node[box, below=of L, align=left] (st) {stage list\\ K: $[S,\hat A_1,\hat A_2,\hat A_3]$ (stem first)\\ E: $[\hat A_1,\hat A_2]$\\ $F_{eeg}$: $[\hat A_1,\hat A_2,C]$ (cross-attention last)};
\node[box, below left=7mm and -16mm of st] (ro) {\texttt{rollout}: product of the matrices; row $t$ gives 400 shares};
\node[box, below right=7mm and -16mm of st] (bn) {\texttt{bin\_matrix}: bins of 4 tokens (80 ms); sources summed, queries averaged};
\node[box, below=of bn] (fl) {\texttt{attention\_flow}: maximum flow from each of the 100 bins to the target bin};
\node[box, below=of fl] (un) {same flow for uniform attention $\tfrac12\frac1n J+\tfrac12 I$};
\node[result, below=of un] (ex) {flow minus the uniform baseline};
\node[result, below=17mm of ro] (rs) {share of each token, the event's own token left out};
\draw[arr] (x)--(tr); \draw[arr] (tr)--(A); \draw[arr] (A)--(L); \draw[arr] (L)--(st);
\draw[arr] (st)--(ro); \draw[arr] (st)--(bn); \draw[arr] (bn)--(fl); \draw[arr] (fl)--(un); \draw[arr] (un)--(ex); \draw[arr] (ro)--(rs);
\end{tikzpicture}
\end{latin}

\caption{پیاده‌سازی \lr{rollout} (چپ) و \lr{attention flow} (راست) برای مسیر زمانی.}
\label{fig:rf-impl}
\end{figure}

\subsection{چرا \lr{attention flow} در این مدل اشباع می‌شود}

مثال دو لایه با توجه یکنواخت روی $n$ موقعیت را در نظر بگیرید (مسیر زمانی \lr{E}). ظرفیت یال از ورودی $j$ به گره میانی $i$ برابر $\tfrac{1}{2n}+\tfrac12\delta_{ij}$ و از $i$ به هدف $t$ برابر $\tfrac{1}{2n}+\tfrac12\delta_{it}$ است. برای $j\ne t$ مسیرها از گره‌های میانی جدا می‌گذرند و هر کدام با کوچک‌ترین یال خود محدود می‌شوند:
\begin{equation}
\mathrm{flow}(j\to t)=\sum_{i=1}^{n}\min\!\Big(\tfrac{1}{2n}+\tfrac12\delta_{ij},\ \tfrac{1}{2n}+\tfrac12\delta_{it}\Big)=n\cdot\tfrac{1}{2n}=\tfrac12,
\end{equation}
و برای خود هدف ($j=t$) نزدیک به ۱. پس در مدلی که هیچ چیز نیاموخته، همه‌ی ورودی‌ها جریان \lr{۰٫۵} دارند و جریان در \lr{token} خود رخداد به ۱ می‌پرد. با \lr{stem} و سه لایه (مسیر \lr{K}) این عدد برای همه‌ی ورودی‌ها حدود ۱ است. وقتی توجه پخش است، تعداد زیاد مسیرهای موازی کوچک جمع می‌شود و جریان به «پخش بودن توجه» حساس است، نه به اینکه اطلاعات از کجا آمده است.

در الکترودها وضع بدتر است. جریان هر الکترود باید از ردیف \lr{pooling} بگذرد که ظرفیت کل آن ۱ است، و ظرفیت خروجی هر الکترود در لایه‌ی اول دست‌کم \lr{۰٫۵} است. با توجه نسبتاً پخش در لایه‌ی دوم، تقریباً همه‌ی الکترودها به یک سقف می‌رسند: در \lr{P1} سهم جریان همه‌ی ۳۲ الکترود بین \lr{۳٫۰۹} و \lr{۳٫۱۵} درصد است (سهم یکنواخت \lr{۳٫۱۲۵} درصد است). شکل‌ها در این حالت نقشه‌ی جریان الکترود را رسم نمی‌کنند و دلیلش را در عنوان می‌نویسند.

\subsection{نتایج}

جدول~\ref{tab:flow} و شکل~\ref{fig:flow_all.png} جریان را در برابر خط پایه برای هر ۱۲ شرکت‌کننده نشان می‌دهند. سه نتیجه:
\begin{enumerate}
\item جریان واقعی همه‌جا \emph{کمتر} از \lr{uniform-attention flow} است (میانگین \lr{−۰٫۱۵} برای \lr{K}، \lr{−۰٫۲۱} برای \lr{E} و \lr{−۰٫۲۴} برای $F_{eeg}$). یعنی توجه آموخته‌شده متمرکزتر از یکنواخت است و راه بسیاری از ورودی‌ها را می‌بندد.
\item برای \lr{K} الگوی زمانی در همه‌ی شرکت‌کننده‌ها یکسان است: ورودی‌های ۰ تا ۳ ثانیه پس از جداشدن کمترین جریان را دارند و ورودی‌های پیش از رخداد بیشترین. این با \lr{rollout} سازگار است: \lr{K} برای سطح فعالیت در لحظه‌ی جداشدن، به حرکت پیش از آن نگاه می‌کند.
\item پراکندگی جریان الکترودها در ۱۰ نفر از ۱۲ کمتر از ۱۰ درصد میانگین است؛ فقط \lr{P2} (۱۲ درصد) از این مرز می‌گذرد. جریان الکترود در این مدل تقریباً هیچ اطلاعات مکانی ندارد.
\end{enumerate}
در \lr{deletion check} (جدول~\ref{tab:interp-del}) جریان برای \lr{K} ضعیف‌ترین روش است (\lr{۰٫۱۱۸} در برابر حدود \lr{۰٫۲۱}) و برای \lr{E} هم‌تراز \lr{plain rollout} و بسیار ضعیف‌تر از \lr{gradient-weighted rollout}. نتیجه‌ی عملی: \lr{attention flow} در این معماری درست محاسبه می‌شود اما برای مکان‌یابی اطلاعات به کار نمی‌آید؛ \lr{gradient-weighted rollout} روش اصلی است و جریان فقط به‌صورت «فاصله از خط پایه‌ی یکنواخت» گزارش می‌شود.

\begin{table}[htbp]
\centering
\small
\caption{\lr{attention flow} در برابر خط پایه‌ی توجه یکنواخت: میانگین «جریان منهای جریانِ توجه یکنواخت» روی همه‌ی بازه‌های زمانی، و پراکندگی جریان بین الکترودها (بیشینه منهای کمینه، تقسیم بر میانگین). منفی یعنی جریان از حالت یکنواخت کمتر است.}
\label{tab:flow}
\begin{tabular}{ccccc}
\toprule
شرکت‌کننده & جریان اضافه \lr{K} & جریان اضافه \lr{E} & جریان اضافه \lr{F\_eeg} & پراکندگی جریان الکترودها (\%) \\
\midrule
\lr{P1} & \lr{−۰٫۱۵۷} & \lr{−۰٫۲۲۸} & \lr{−۰٫۲۴۶} & \lr{۱٫۷} \\
\lr{P2} & \lr{−۰٫۲۴۰} & \lr{−۰٫۲۰۳} & \lr{−۰٫۲۰۴} & \lr{۱۱٫۵} \\
\lr{P3} & \lr{−۰٫۱۱۶} & \lr{−۰٫۲۴۲} & \lr{−۰٫۲۰۲} & \lr{۴٫۵} \\
\lr{P4} & \lr{−۰٫۱۳۲} & \lr{−۰٫۲۱۶} & \lr{−۰٫۲۱۷} & \lr{۱٫۶} \\
\lr{P5} & \lr{−۰٫۱۴۱} & \lr{−۰٫۱۹۵} & \lr{−۰٫۱۸۶} & \lr{۶٫۳} \\
\lr{P6} & \lr{−۰٫۱۸۳} & \lr{−۰٫۱۷۹} & \lr{−۰٫۱۹۱} & \lr{۵٫۶} \\
\lr{P7} & \lr{−۰٫۱۸۸} & \lr{−۰٫۱۶۸} & \lr{−۰٫۲۱۶} & \lr{۱٫۴} \\
\lr{P8} & \lr{−۰٫۰۷۶} & \lr{−۰٫۱۴۸} & \lr{−۰٫۲۷۵} & \lr{۸٫۸} \\
\lr{P9} & \lr{−۰٫۱۲۱} & \lr{−۰٫۲۷۰} & \lr{−۰٫۲۴۳} & \lr{۱٫۳} \\
\lr{P10} & \lr{−۰٫۲۰۶} & \lr{−۰٫۱۵۶} & \lr{−۰٫۳۹۱} & \lr{۰٫۹} \\
\lr{P11} & \lr{−۰٫۰۷۷} & \lr{−۰٫۲۲۳} & \lr{−۰٫۲۴۰} & \lr{۲٫۴} \\
\lr{P12} & \lr{−۰٫۱۲۸} & \lr{−۰٫۳۰۰} & \lr{−۰٫۲۵۸} & \lr{۲٫۸} \\
\bottomrule
\end{tabular}
\end{table}

\fig{flow_all.png}{1.0}{\lr{attention flow} منهای \lr{uniform-attention flow} برای \lr{K} (چپ) و \lr{E} (وسط) حول جداشدن جسم، هر خط یک شرکت‌کننده؛ بازه‌ی \lr{token} خود رخداد حذف شده است. راست: پراکندگی جریان الکترودها برای \lr{E}؛ خط‌چین مرز ۱۰ درصد که زیر آن نقشه‌ی جریان رسم نمی‌شود.}
